%%%%%%%%%%%%%%%%%
% This is an sample CV template created using altacv.cls
% (v1.6, 21 May 2021) written by LianTze Lim (liantze@gmail.com). Now compiles with pdfLaTeX, XeLaTeX and LuaLaTeX.
%
%% It may be distributed and/or modified under the
%% conditions of the LaTeX Project Public License, either version 1.3
%% of this license or (at your option) any later version.
%% The latest version of this license is in
%%    http://www.latex-project.org/lppl.txt
%% and version 1.3 or later is part of all distributions of LaTeX
%% version 2003/12/01 or later.t
%%%%%%%%%%%%%%%%

%% Use the "normalphoto" option if you want a normal photo instead of cropped to a circle
% \documentclass[10pt,a4paper,normalphoto]{altacv}

\documentclass[10pt,a4paper,ragged2e,withhyper]{altacv}
%% AltaCV uses the fontawesome5 and packages.
%% See http://texdoc.net/pkg/fontawesome5 for full list of symbols.

% Change the page layout if you need to
\geometry{left=1.25cm,right=1.25cm,top=1.5cm,bottom=1.5cm,columnsep=1.2cm}

% The paracol package lets you typeset columns of text in parallel
\usepackage{paracol}
\usepackage{hyperref}
\usepackage{fontawesome5}


% Change the font if you want to, depending on whether
% you're using pdflatex or xelatex/lualatex
\ifxetexorluatex
  % If using xelatex or lualatex:
  \setmainfont{Roboto Slab}
  \setsansfont{Lato}
  \renewcommand{\familydefault}{\sfdefault}
\else
  % If using pdflatex:
  \usepackage[rm]{roboto}
  \usepackage[defaultsans]{lato}
  % \usepackage{sourcesanspro}
  \renewcommand{\familydefault}{\sfdefault}
\fi

% Change the colours if you want to
\definecolor{SlateGrey}{HTML}{2E2E2E}
\definecolor{LightGrey}{HTML}{666666}
\definecolor{DarkPastelRed}{HTML}{450808}
\definecolor{PastelRed}{HTML}{8F0D0D}
\definecolor{GoldenEarth}{HTML}{E7D192}
% \colorlet{name}{black}
% \colorlet{tagline}{NewBlue}
% \colorlet{heading}{DarkPastelRed}
% % \colorlet{headingrule}{BDE5E1}
% \colorlet{subheading}{NewBlue}
% \colorlet{accent}{NewBlue}
% \colorlet{emphasis}{SlateGrey} % should be fine
% \colorlet{body}{LightGrey}

\definecolor{LightBlue}{HTML}{6BBED4}
\definecolor{MidBlue}{HTML}{137dc9}
\definecolor{DarkerBlue}{HTML}{2864cc}

\colorlet{name}{black}
\colorlet{tagline}{MidBlue}
\colorlet{heading}{DarkerBlue}
\colorlet{headingrule}{DarkerBlue}
\colorlet{subheading}{MidBlue}
\colorlet{accent}{MidBlue}
\colorlet{emphasis}{SlateGrey} % should be fine
\colorlet{body}{LightGrey}

% Change some fonts, if necessary
\renewcommand{\namefont}{\Huge\rmfamily\bfseries}
\renewcommand{\personalinfofont}{\footnotesize}
\renewcommand{\cvsectionfont}{\LARGE\rmfamily\bfseries}
\renewcommand{\cvsubsectionfont}{\large\bfseries}

% Change the bullets for itemize and rating marker
% for \cvskill if you want to
\renewcommand{\itemmarker}{{\small\textbullet}}
\renewcommand{\ratingmarker}{\faCircle}
%% Use (and optionally edit if necessary) this .tex if you
%% want to use an author-year reference style like APA(6)
%% for your publication list
\input{pubs-authoryear}

%% Use (and optionally edit if necessary) this .tex if you
%% want an originally numerical reference style like IEEE
%% for your publication list
% \input{pubs-num}

%% sample.bib contains your publications
\addbibresource{sample.bib}
\begin{document}
\name{Meet Vora}
% TODO [TAGLINE] — Update tagline once AI experience is stronger.
%   Current tagline emphasizes Meta/reliability. Once flagship project or deeper AI work
%   is done, consider something like: "Software Engineer · Systems & AI · ..."
\tagline{Software Engineer · 3+ Years at Meta · Building Reliable Systems at Scale}
%% You can add multiple photos on the left or right
% \photoR{2.8cm}{Globe_High}
% \photoL{2.5cm}{Yacht_High,Suitcase_High}

\personalinfo{%
  % Not all of these are required!
  \email{meetrajivvora@gmail.com}
  \phone{(209) 891-7885}
%   \mailaddress{Åddrésş, Street, 00000 Cóuntry}
%   \location{Location, COUNTRY}
%   \homepage{www.homepage.com}
%   \twitter{@twitterhandle}
  \linkedin{linkedin.com/in/mrvora}
  \github{github.com/meet-vora}
%   \orcid{0000-0000-0000-0000}
  %% You can add your own arbitrary detail with
  %% \printinfo{symbol}{detail}[optional hyperlink prefix]
  % \printinfo{\faPaw}{Hey ho!}[https://example.com/]
  %% Or you can declare your own field with
  %% \NewInfoFiled{fieldname}{symbol}[optional hyperlink prefix] and use it:
  % \NewInfoField{gitlab}{\faGitlab}[https://gitlab.com/]
  % \gitlab{your_id}
  %%
  %% For services and platforms like Mastodon where there isn't a
  %% straightforward relation between the user ID/nickname and the hyperlink,
  %% you can use \printinfo directly e.g.
  % \printinfo{\faMastodon}{@username@instace}[https://instance.url/@username]
  %% But if you absolutely want to create new dedicated info fields for
  %% such platforms, then use \NewInfoField* with a star:
  % \NewInfoField*{mastodon}{\faMastodon}
  %% then you can use \mastodon, with TWO arguments where the 2nd argument is
  %% the full hyperlink.
  % \mastodon{@username@instance}{https://instance.url/@username}
}

\makecvheader
%% Depending on your tastes, you may want to make fonts of itemize environments slightly smaller
% \AtBeginEnvironment{itemize}{\small}

%% Set the left/right column width ratio to 6:4.
\columnratio{0.6}

% Start a 2-column paracol. Both the left and right columns will automatically
% break across pages if things get too long.
\begin{paracol}{2}
\cvsection{Experience}

\cvevent{Software Engineer (IC4)}{Meta}{September 2022 - March 2026}
\begin{itemize}
\item Delivered critical compliance platform features end-to-end for Project Compass (1K+ MAU, \textasciitilde\$14M in savings), launching 9 months ahead of schedule across multiple domain rollouts
\item Sole engineer on compliance review tooling --- reduced end-to-end review SLA (135 \textrightarrow{} 43 days), cut \textasciitilde\$350k/yr in legal review costs by independently delivering all milestones ahead of schedule
\item Reduced API query volume by 99.9\% (27M \textrightarrow{} 20K hits) and boosted query performance by 33\%. Rebuilt compliance PowerSearch system from scratch with reliable operators and restructured search fields
\item Integrated ML-generated taxonomy recommendations into compliance tooling, replacing manual classification workflows
\item Demoed technical work to Director and VP leadership. Reviewed 400+ PRs, mentored new hires, and independently drove alignment across 10+ cross-functional partners
\end{itemize}

\medskip

\cvsection{AI \& Automation}

\begin{itemize}
\item Built an AI-first developer productivity system at Meta: authored 5 custom rules, 12+ skills/commands, and 10+ hooks, connected to 4 autonomous agents running 24/7 on a dedicated server
% TODO [AI PROJECT BULLET] — Add a second strong bullet here once a real AI project is complete.
%   This could be the flagship Claude API project or another meaningful AI/ML project.
%   The AI section currently has only 1 real bullet + 1 vague one — needs at least 2 strong ones.
% TODO [FLAGSHIP AI PROJECT] — Build a standout project using Claude's API.
%   Ideas: CLI code review tool, safety eval harness, RAG app over a specific domain.
%   Put on GitHub with clean README. Replace this placeholder with 1-line description + link.
%   This is the single highest-signal item for Anthropic. See Prep Plan Phase 4, Item 3.
% \item [PLACEHOLDER: Flagship Claude API project — e.g. "Built [tool name], a [description] using Claude's API (github.com/meet-vora/[repo])"]
%
% TODO [OPEN SOURCE CONTRIBUTION] — Optional but high-impact.
%   Targets: anthropic-sdk-python, anthropic-cookbook, LangChain, or any Python OSS project.
%   Even a small, well-crafted PR counts. See Prep Plan Phase 4, Item 2.
% \item [PLACEHOLDER: Open-source contribution — e.g. "Contributed [feature/fix] to [repo]"]
%
% TODO [PERSONAL AI WORKFLOW] — Flesh out once home setup is more concrete.
%   Currently too vague for Anthropic. To make this a real bullet, build 1-2 of:
%   - Scheduled tasks + custom skills in Cowork (turns "I use AI" into "I built AI automation")
%   - RAG pipeline over Obsidian vault (semantic search across notes)
%   - Automated knowledge ingestion (articles → summarized → tagged → filed)
%   - Personal CRM with AI-powered follow-up nudges
%   Once built, replace with something like:
%   "Built personal AI automation layer over Obsidian: [N] scheduled tasks, [N] custom skills,
%    and a [specific feature] for [specific outcome]"
% \item Building personal AI-augmented workflow for knowledge management and task automation
\end{itemize}

\medskip

\cvsection{Prior Experience}

\cvevent{Backend SWE Intern}{Zuora}{January 2022 - April 2022}
\begin{itemize}
\item Built Spring API with 7 endpoints for Kafka-Connect and Debezium connectors, pipelining data from MySQL \& Oracle databases to Kafka for company-wide services
\end{itemize}

\divider

\cvevent{Fullstack SWE Intern}{Checkbook}{March 2021 - August 2021}
\begin{itemize}
\item Implemented features end-to-end across 4+ repositories. Rebuilt React site structure for internationalization, reducing code duplication by 60\%
\end{itemize}

% \divider

% \cvevent{Gitlet}{Version Control System}{April 2020 - May 2020}
% \begin{itemize}
%     \item Git-like version control system that tracks files in a local repository
%     \item Commands to track files include init, add, commit, rm, log, branch, checkout, and more!
%     \item Wrote 1000+ lines of code in Java and spent 20+ hours to implement
% \end{itemize}

% \medskip

% use ONLY \newpage if you want to force a page break for
% ONLY the current column
% \newpage

% \cvsection{Publications}

% \nocite{*}

% \printbibliography[heading=pubtype,title={\printinfo{\faBook}{Books}},type=book]

% \divider

% \printbibliography[heading=pubtype,title={\printinfo{\faFile*[regular]}{Journal Articles}},type=article]

% \divider

% \printbibliography[heading=pubtype,title={\printinfo{\faUsers}{Conference Proceedings}},type=inproceedings]

%% Switch to the right column. This will now automatically move to the second
%% page if the content is too long.
\switchcolumn

% \cvsection{My Life Philosophy}

% \begin{quote}
% ``Something smart or heartfelt, preferably in one sentence.''
% \end{quote}

\cvsection{Summary}
\begin{itemize}
    \item 3+ years building reliable backend systems at Meta scale
    \item End-to-end project ownership across compliance, data, and platform tooling
    \item Proven track record of quantifiable impact: \$14M+ savings, 68\% SLA reduction
    \item Applying lessons from privacy-critical systems at Meta to building trustworthy AI
\end{itemize}

% TODO [BLOG POST / WRITING] — Anthropic explicitly values this. See Prep Plan Phase 4, Item 1.
%   Write about: a system you designed (keep non-confidential), system design analysis,
%   or your perspective on AI safety. Publish on Medium or personal site.
%   Uncomment below and add link once published.
% \cvsection{Writing}
% \begin{itemize}
%     \item [PLACEHOLDER: "Title of Blog Post" — medium.com/@meetvora/...]
% \end{itemize}

\cvsection{Languages}

\cvtag{Python}
\cvtag{Hack/PHP}
\cvtag{TypeScript}
\cvtag{Java}
\cvtag{JavaScript}
\cvtag{SQL}
\cvtag{C}

\medskip

\cvsection{Frameworks/Tools}

\textbf{Backend:}
\cvtag{GraphQL}
\cvtag{REST APIs}
\cvtag{Spring Boot}
\cvtag{Flask}
\cvtag{PostgreSQL}
\cvtag{Kafka}
\divider

\textbf{Frontend:}
\cvtag{React}
\divider

% TODO [AI/ML SKILLS] — Uncomment once you have hands-on experience to back these up.
%   After flagship project: add Claude API, prompt engineering
%   After RAG project: add LangChain or LlamaIndex
%   After OSS contribution: add relevant framework
% \textbf{AI/ML:}
% \cvtag{Claude API}
% \cvtag{Prompt Engineering}
% \cvtag{RAG}
% \divider

\textbf{Tools:}
\cvtag{Git}
\cvtag{Linux}
\cvtag{Docker}
\cvtag{CI/CD}
% \cvtag{Big Data}

% \cvsection{Most Proud of}


% \divider


% \divider



% \cvsection{Languages}

% \cvskill{English}{5}
% \divider

% \cvskill{Spanish}{4}
% \divider

% \cvskill{German}{3.5} %% Supports X.5 values.

%% Yeah I didn't spend too much time making all the
%% spacing consistent... sorry. Use \smallskip, \medskip,
%% \bigskip, \vspace etc to make adjustments.
\medskip

\cvsection{Education}

\cvgrad{\textbf{University of California, Berkeley}}{ Computer Science \kern 5pt}{August 2019 - May 2022}

\begin{itemize}
\item Graduated one full year early
\end{itemize}

\medskip

\cvsection{Hobbies}
\cvtag{Fantasy Novels}
\cvtag{Tennis}
\cvtag{Backpacking}
\cvtag{Board Games}

% \Large\color{LightBlue}\faBook
% \hspace{18pt}
% \Large\color{LightBlue}\faGamepad
% \hspace{18pt}
% \Large\color{LightBlue}\faUsers
% \hspace{18pt}
% \Large\color{LightBlue}\faHeadphones
% \hspace{18pt}
% \Large\color{LightBlue}\faBasketballBall
% \hspace{20pt}

% \divider



\end{paracol}


\end{document}
